To our knowledge, this is the first paper to (1)~demonstrate the computational inefficiencies in the GenASM algorithm, (2)~address them with three improvements in our new \toolname algorithm, (3)~rigorously demonstrate the computational benefits of \toolname over GenASM for CPU, GPU, ASIC implementations, and (4)~rigorously analyze the accuracy of GenASM and \toolname under multiple different configurations.

We have already extensively compared to WFA~\citep{wfa_algorithm}, KSW2~\citep{ksw2_a, ksw2_b}, Edlib~\citep{edlib}, CUDASW++3.0~\citep{liu2013cudasw++}, and Darwin-GPU~\citep{ahmed2020darwin_gpu}. Several other works accelerate sequence alignment: NVBIO~\citep{nvbio} is a multi-purpose library for accelerating bioinformatics applications using GPUs, but is no longer maintained. Gasal2~\citep{gasal2} is a recent GPU aligner limited to short reads. CUDAlign4.0 \citep{de2016cudalign} can efficiently align a single pair of extremely long (chromosome-sized) sequences, with use cases such as whole genome alignment. Adept~\citep{awan2020adept} is a recent GPU aligner for short and long reads but does not support traceback, i.e. only reports the alignment score.

The Darwin accelerator~\citep{turakhia2019darwin} implements a Smith-Waterman-Gotoh accelerator for long reads using a similar greedy strategy to GenASM called \defn{tiling}. We have compared \toolname to the GPU implementation of this algorithm, Darwin-GPU. GenASM, \toolname, and Darwin demonstrate the significant benefits of greedy algorithms, based on which there are at least two interesting future directions to explore. First, a suitability study of different algorithms to greedy heuristics, such as Myers' bitvector algorithm~\citep{myers_algorithm_1999}, Hyyrö's banded bitvector algorithm~\citep{hyyro2003bit} or the recently proposed wavefront algorithm~\citep{wfa_algorithm}. Second, an exploration of the effectiveness of our algorithmic improvements for other implementations of greedy windowing or tiling, like Darwin. We believe the DENT improvement can be applied directly to Darwin.

We have demonstrated the computational benefits of \toolname over a variety of state-of-the-art baselines for both commodity hardware (i.e., CPUs and GPUs) and custom hardware (i.e., ASICs). We have demonstrated the accuracy of \toolname for multiple datasets. We conclude that \toolname\ has clear benefits across a wide range of computing platforms.